%\chapter{Reference guide}

%
\section{Parallel processes with IPC}
From Astra--6.2.1 on, it is possible to run several processes on independent memory segments, to speed up computations in cases where
each radial grid point is independently considered.

As an example, suppose to run a transport simulation of electron temperature, for which the energy diffusivity is computed with a sophisticated model that requires some time for a single radial point computation.
Said model is local, i.e. each radial grid point is computed independently of the others.
In this case, the IPC paradigm is of help, since many radial grid points can be then run simultaneously which greatly speeds up the overall simulation.
%
\subsection{Calling IPC from the model file}
Suppose to have a subroutine {\it subp.f} which calls a physical model that computes at generic radii index $j$. \\
In the model file the IPC call to {\it subp.f} appears as a sequence of calls:

xpr/subp(j1,j2)\&::::;	\\
where $j1,j2$ are indexes of radial positions. As an example, with 121 radial grid points:

xpr/subp(1,40)\&::::;	

xpr/subp(41,81)\&::::;	

xpr/subp(82,121)\&::::;	\\
which means that the subp module is sent on three independent processes, the first of which computes radii from 1 to 40, the second from 41 to 81, and the third from 82 to 121. In this way, the computational time for 121 grid points is 1/3 than the usual computational time that would require the call 

subp::::;\\
which solves all radii in one single process.
%
\subsection{Setting up the subroutines to run the IPC processes}
Running IPC requires some effort from the part of the user in setting up a bunch
of subroutines that take care of appropriate variable assignment and calls. This
is unfortunately not as straightforward as calling subroutines in the standard 
way.

The detailed explanation on how to set up these additional routines and the scheme of the IPC process is presented in the attached paper {\it IPC.pdf}.
%
\subsection{Calling subroutines in order}
In general, in all Astra versions, the order in which commands enter in the model file, is maintained in the compilation. In addition to this, from Astra--6.2.1 onwards, there is the possibility to explicitly tell the code to give higher or lower priority to specific subroutines.

Suppose to have five generic subroutines that here we call SUB1,SUB2,SUB3,SUB4,SUB5, where the number is arbitrary, does not represent priority.
If in the model file they appear for example as:\\
SUB1::::;\\
SUB2::::;\\
SUB3::::;\\
SUB4::::;\\
Means that they will be performed in the order in which they appear: 1,2,3,4.
Now, by writing for example:\\
SUB1::::;\\
SUB2$<$::::;\\
SUB3::::;\\
SUB4::::;\\
The order would be: 2,1,3,4. SUB2 has acquired highest priority with the symbol $<$.
The following instead:\\
SUB1::::;\\
SUB2$<$::::;\\
SUB3$<$::::;\\
SUB4::::;\\
gives the order: 2,3,1,4.
In this other example:\\
SUB1$>$::::;\\ 
SUB2::::;\\
SUB3::::;\\
SUB4::::;\\
the order is: 2,3,4,1. So SUB1 has lowest priority with the symbol $>$.
In this example:\\
SUB1$>$::::;\\
SUB2$>$::::;\\
SUB3::::;\\
SUB4::::;\\
the order is then 3,4,1,2.



%

\clearpage
